\section{ĐỀ SỐ \thesection } %star6
\setcounter{bt}{0}

%%Câu 1.
\begin{bt} Thực hiện phép tính:
		\begin{listEX}[2]	
			\item $354 + 307 + 46 + 193$;
			\item $375 : \left\{ 32 - \left[ 4 + (5 \cdot 3^2 - 42) \right] \right\} - 14$;
			\item $3^5 \cdot 108 - 3^5 \cdot 107 + 9^2 \cdot 3$;
			\item $4^3 \cdot 5 - 200 : 5^2 + 2023^0$.
		\end{listEX}
	\loigiai{a) $900$ \hfill b) $1$ \hfill c) $486$ \hfill d) $313$.}
\end{bt}


%%Câu 2.
\begin{bt} Tìm số nguyên $x$, biết:
		\begin{listEX}[2]
		 \item $x - 48 : 16 = 37 \cdot 2$;
		\item $6(x + 2^3) + 40 = 100$;
		\item $3(x + 10) - 15 = 19+x$;
		\item $2^{x+1}+5^2\cdot 2^x=2^3\cdot 3^3$.
		\end{listEX}
	\loigiai{a) $x = 77$ \hfill b) $x = 2$ \hfill c) $x = 2$ \hfill d) $x = 3$.
	}
\end{bt}


%%Câu 5.
\begin{bt} Một rạp xiếc khuyến mãi giảm giá $5\%$ cho vé người lớn và $15\%$ cho vé trẻ em. Biết giá niêm yết của vé trẻ em là $60\,000$ đồng.
	\begin{listEX}
		\item Tính giá vé trẻ em sau khi được giảm giá; 
		\item Biết giá sau khi giảm của vé người lớn là $123\,500$ đồng. Hãy tính giá niêm yết của vé người lớn
	\end{listEX}
	\loigiai{
		a) $51\,000$ đồng \hfill b) $130\,000$ đồng
		}
\end{bt}


%%Câu 3.
\begin{bt} Ba ô tô cùng khởi hành một lúc từ một bến. Thời gian cả đi lẫn về của xe thứ nhất là $40$ phút, của xe thứ hai là $50$ phút, của xe thứ ba là $30$ phút. Khi trở về bến, mỗi xe đều nghỉ $10$ phút rồi tiếp tục chạy. Hỏi sau ít nhất bao lâu:
	\begin{listEX}[2]
	 \item Xe thứ nhất và xe thứ hai cùng rời bến?
	\item Xe thứ hai và xe thứ ba cùng rời bến?
	\item Cả ba xe cùng rời bến?
	\end{listEX}
	\loigiai{ 
	a) $5$ giờ \hfill b) $2$ giờ \hfill c) $10$ giờ	}
\end{bt}


%%Câu 4.
\begin{bt} Để phòng dịch Covid-19, xã ABC thành lập các đội phản ứng nhanh gồm bác sĩ và điều dưỡng viên. Ban đầu có $45$ bác sĩ và $55$ điều dưỡng viên đăng kí. Khi chia đội thực tế có thêm $3$ bác sĩ đăng kí tham gia và có $1$ điều dưỡng viên bị nhiễm Covid-19 nên không tham gia. Hỏi có thể thành lập được nhiều nhất bao nhiêu đội phản ứng nhanh? Khi đó mỗi đội có bao nhiêu bác sĩ và bao nhiêu điều dưỡng viên?
	\loigiai{ $6$ đội; mỗi đội $8$ bác sĩ và $9$ điều dưỡng viên.}
\end{bt}



\begin{bt} Nhà thi đấu Phú Thọ có sàn hình chữ nhật, chiều dài $130$ m và chiều rộng $105$ m. 
	\begin{listEX}
		\item Tính diện tích sàn xây dựng của nhà thi đấu Phú Thọ;
		\item Người ta lát sàn bằng gạch hình vuông cạnh $60$ cm. Tính số thùng gạch cần để lát hết sàn, biết hao hụt khi thi công là $5\%$ số gạch và mỗi thùng có $50$ viên gạch.
	\end{listEX}
	\loigiai{a) $13\,650$ m$^2$ \hfill b) $797$ thùng
	}
\end{bt}



\begin{bt}Kết quả điều tra về loại quả ưa thích nhất của một số bạn trong lớp được ghi như bảng 
\begin{center}
		\begin{tabular}{|c|c|c|c|}

		\hline
		Cam & Ổi & Chuối & Xoài \\ \hline
		Khế & Cam & Ổi & Khế \\ \hline
		Xoài & Ổi & Chuối & Xoài \\ \hline
		Xoài & Cam & Khế & Cam \\ \hline
		Ổi & Khế & Xoài & Chuối \\ \hline
	\end{tabular}
\end{center}
	\begin{listEX}
		\item Có bao nhiêu bạn tham gia trả lời?
		\item Hãy lập bảng thống kê và cho biết loại quả nào được yêu thích nhất; loại quả đó chiếm bao nhiêu phần trăm so với tổng số bạn
	\end{listEX}
	\loigiai{
		\begin{listEX}
			\item   $20$ bạn tham gia
			\item Thống kê tần số và tỉ lệ:
			
			\begin{center}
				\begin{tabular}{|c|c|c|}
					\hline
					\textbf{Loại quả} & \textbf{Số bạn} & \textbf{Tỉ lệ} \\ \hline
					Cam & $4$ & $20\%$ \\ \hline
					Ổi & $4$ & $20\%$ \\ \hline
					Chuối & $3$ & $15\%$ \\ \hline
					Xoài & $5$ & $25\%$ \\ \hline
					Khế & $4$ & $20\%$ \\ \hline
					\textbf{Tổng} & $20$ & $100\%$ \\ \hline
				\end{tabular}
			\end{center}
			
			Vậy quả được yêu thích nhất là Xoài, chiếm $25\%$ tổng số bạn
	\end{listEX}}
\end{bt}

\begin{bt} 
	Tìm số tự nhiên có năm chữ số, biết rằng nếu viết thêm chữ số $2$ vào đằng sau số đó thì được số lớn gấp ba lần số có được bằng cách viết thêm chữ số $2$ vào đằng trước số ban đầu.
	\loigiai{
		Gọi số phải tìm là $\overline{abcde} = x$ \quad ($a,b,c,d,e \in \mathbb{N}$; $a \ne 0$; $a,b,c,d,e \le 9$).
		Khi đó:
		\begin{align*}
			\overline{abcde2} &= 3 \cdot \overline{2abcde} \\
			10x + 2 &= 3\cdot(200000 + x) \\
			x &= 85714.
		\end{align*}
		
	}
\end{bt}
